% =====================================================================
%  2.4 Ulikheter — teori og oppgaver
% =====================================================================

\del{1}{Lineære ulikheter}

\begin{ordliste}
ulikhet & To uttrykk med et ulikhetstegn mellom. Løsningen er ofte mange tall. & \mittsprak\\
ulikhetstegn & Tegnene $<$, $>$, $\le$ og $\ge$. & \mittsprak\\
tallinje & En linje der hvert tall har sin plass. Vi viser løsningen på den. & \mittsprak\\
intervall & Alle tallene mellom to endepunkter, for eksempel $\langle 1, 4]$. & \mittsprak\\
åpent og lukket intervall & Åpent: endepunktet er ikke med ($\langle$ eller $\rangle$). Lukket: endepunktet er med ($[$ eller $]$). & \mittsprak\\
\end{ordliste}

\begin{mustkunne}
\begin{minipage}[t]{0.46\linewidth}
\textbf{Regler.} Vi løser en ulikhet som en likning, men:
\begin{itemize}[leftmargin=1.2em,itemsep=1pt,topsep=2pt]
  \item Legg til eller trekk fra det samme på begge sider.
  \item Gang eller del på det samme \emph{positive} tallet.
  \item Gang eller del på et \emph{negativt} tall: \textbf{snu} ulikhetstegnet.
\end{itemize}
\end{minipage}\hfill
\begin{minipage}[t]{0.50\linewidth}
\textbf{Hvorfor snur vi tegnet?} Vi vet at $2 < 3$. Gang med $-1$. Tallene speiles om 0, og rekkefølgen blir motsatt: $-2 > -3$.

\medskip
\begin{tikzpicture}[x=0.6cm,font=\small]
  \draw[-{Stealth[length=2mm]}] (-4.6,0) -- (4.6,0);
  \foreach \i in {-4,...,4} \draw (\i,0.07) -- (\i,-0.07) node[below,font=\scriptsize,inner sep=2pt] {$\i$};
  \fill[hoved] (2,0) circle (2.5pt);
  \fill[hoved] (3,0) circle (2.5pt);
  \fill[niv3] (-2,0) circle (2.5pt);
  \fill[niv3] (-3,0) circle (2.5pt);
  \draw[->,graa] (2,0.2) to[bend right=60] (-2,0.2);
  \draw[->,graa] (3,0.2) to[bend right=60] (-3,0.2);
\end{tikzpicture}
\end{minipage}

\medskip
\textbf{Tre måter å svare på.} Åpen ring og $\langle$ eller $\rangle$: tallet er ikke med. Fylt ring og $[$ eller $]$: tallet er med.
\renewcommand{\arraystretch}{1.3}

\smallskip
\begin{tabular}{@{}l l l@{}}
\textbf{Ulikhet} & \textbf{Tallinje} & \textbf{Intervall}\\
$x > 2$ & \begin{tikzpicture}[x=0.45cm,font=\small]
  \draw[-{Stealth[length=2mm]}] (-2.5,0) -- (5.5,0);
  \draw (-2,0.07) -- (-2,-0.07) node[below,font=\scriptsize,inner sep=2pt] {$-2$};
  \draw (-1,0.07) -- (-1,-0.07) node[below,font=\scriptsize,inner sep=2pt] {$-1$};
  \draw (0,0.07) -- (0,-0.07) node[below,font=\scriptsize,inner sep=2pt] {$0$};
  \draw (1,0.07) -- (1,-0.07) node[below,font=\scriptsize,inner sep=2pt] {$1$};
  \draw (2,0.07) -- (2,-0.07) node[below,font=\scriptsize,inner sep=2pt] {$2$};
  \draw (3,0.07) -- (3,-0.07) node[below,font=\scriptsize,inner sep=2pt] {$3$};
  \draw (4,0.07) -- (4,-0.07) node[below,font=\scriptsize,inner sep=2pt] {$4$};
  \draw (5,0.07) -- (5,-0.07) node[below,font=\scriptsize,inner sep=2pt] {$5$};
  \draw[hoved,line width=2.2pt] (2.0,0) -- (5.5,0);
  \filldraw[hoved,fill=white,line width=0.9pt] (2.0,0) circle (2.6pt);
\end{tikzpicture} & $\langle 2, \rightarrow\rangle$\\
$x \le 3$ & \begin{tikzpicture}[x=0.45cm,font=\small]
  \draw[-{Stealth[length=2mm]}] (-2.5,0) -- (5.5,0);
  \draw (-2,0.07) -- (-2,-0.07) node[below,font=\scriptsize,inner sep=2pt] {$-2$};
  \draw (-1,0.07) -- (-1,-0.07) node[below,font=\scriptsize,inner sep=2pt] {$-1$};
  \draw (0,0.07) -- (0,-0.07) node[below,font=\scriptsize,inner sep=2pt] {$0$};
  \draw (1,0.07) -- (1,-0.07) node[below,font=\scriptsize,inner sep=2pt] {$1$};
  \draw (2,0.07) -- (2,-0.07) node[below,font=\scriptsize,inner sep=2pt] {$2$};
  \draw (3,0.07) -- (3,-0.07) node[below,font=\scriptsize,inner sep=2pt] {$3$};
  \draw (4,0.07) -- (4,-0.07) node[below,font=\scriptsize,inner sep=2pt] {$4$};
  \draw (5,0.07) -- (5,-0.07) node[below,font=\scriptsize,inner sep=2pt] {$5$};
  \draw[hoved,line width=2.2pt] (-2.5,0) -- (3.0,0);
  \filldraw[hoved,fill=hoved,line width=0.9pt] (3.0,0) circle (2.6pt);
\end{tikzpicture} & $\langle\leftarrow, 3]$\\
$-1 \le x < 4$ & \begin{tikzpicture}[x=0.45cm,font=\small]
  \draw[-{Stealth[length=2mm]}] (-2.5,0) -- (5.5,0);
  \draw (-2,0.07) -- (-2,-0.07) node[below,font=\scriptsize,inner sep=2pt] {$-2$};
  \draw (-1,0.07) -- (-1,-0.07) node[below,font=\scriptsize,inner sep=2pt] {$-1$};
  \draw (0,0.07) -- (0,-0.07) node[below,font=\scriptsize,inner sep=2pt] {$0$};
  \draw (1,0.07) -- (1,-0.07) node[below,font=\scriptsize,inner sep=2pt] {$1$};
  \draw (2,0.07) -- (2,-0.07) node[below,font=\scriptsize,inner sep=2pt] {$2$};
  \draw (3,0.07) -- (3,-0.07) node[below,font=\scriptsize,inner sep=2pt] {$3$};
  \draw (4,0.07) -- (4,-0.07) node[below,font=\scriptsize,inner sep=2pt] {$4$};
  \draw (5,0.07) -- (5,-0.07) node[below,font=\scriptsize,inner sep=2pt] {$5$};
  \draw[hoved,line width=2.2pt] (-1.0,0) -- (4.0,0);
  \filldraw[hoved,fill=hoved,line width=0.9pt] (-1.0,0) circle (2.6pt);
  \filldraw[hoved,fill=white,line width=0.9pt] (4.0,0) circle (2.6pt);
\end{tikzpicture} & $[-1, 4\rangle$\\
\end{tabular}
\end{mustkunne}

\begin{eksempel}[1 — del på et negativt tall]
Løs $5 - 3x \ge 11$.

\begin{stegtabell}
$-3x \ge 6$ & Trekk fra 5 på begge sider.\\
$x \le -2$ & Del på $-3$. Tallet er negativt, så vi snur tegnet.\\
$\langle\leftarrow, -2]$ & Intervall. $-2$ er med, så klammen er $]$.\\
Kontroll: & $x = -3$ gir $5 + 9 = 14 \ge 11$, som stemmer. $x = 0$ gir $5 \ge 11$, som er usant.\\
\end{stegtabell}
\end{eksempel}

\begin{eksempel}[2 — dobbel ulikhet]
Løs $-1 < 2x + 3 \le 7$. Vis svaret på tallinjen.

\begin{stegtabell}
$-4 < 2x \le 4$ & Trekk fra 3 i alle tre deler.\\
$-2 < x \le 2$ & Del på 2 i alle deler. 2 er positivt, så tegnene beholdes.\\
$\langle -2, 2]$ & Intervall. $-2$ er ikke med ($\langle$), 2 er med ($]$).\\
\end{stegtabell}

\smallskip
\begin{center}
\begin{tikzpicture}[x=0.7cm,font=\small]
  \draw[-{Stealth[length=2mm]}] (-4.5,0) -- (4.5,0);
  \draw (-4,0.07) -- (-4,-0.07) node[below,font=\scriptsize,inner sep=2pt] {$-4$};
  \draw (-3,0.07) -- (-3,-0.07) node[below,font=\scriptsize,inner sep=2pt] {$-3$};
  \draw (-2,0.07) -- (-2,-0.07) node[below,font=\scriptsize,inner sep=2pt] {$-2$};
  \draw (-1,0.07) -- (-1,-0.07) node[below,font=\scriptsize,inner sep=2pt] {$-1$};
  \draw (0,0.07) -- (0,-0.07) node[below,font=\scriptsize,inner sep=2pt] {$0$};
  \draw (1,0.07) -- (1,-0.07) node[below,font=\scriptsize,inner sep=2pt] {$1$};
  \draw (2,0.07) -- (2,-0.07) node[below,font=\scriptsize,inner sep=2pt] {$2$};
  \draw (3,0.07) -- (3,-0.07) node[below,font=\scriptsize,inner sep=2pt] {$3$};
  \draw (4,0.07) -- (4,-0.07) node[below,font=\scriptsize,inner sep=2pt] {$4$};
  \draw[hoved,line width=2.2pt] (-2.0,0) -- (2.0,0);
  \filldraw[hoved,fill=white,line width=0.9pt] (-2.0,0) circle (2.6pt);
  \filldraw[hoved,fill=hoved,line width=0.9pt] (2.0,0) circle (2.6pt);
\end{tikzpicture}
\end{center}
\end{eksempel}

\begin{vanligfeil}
\begin{itemize}[leftmargin=1.2em,itemsep=1pt,topsep=0pt]
  \item $-3x < 12$, så $x < -4$ er feil. Vi deler på $-3$, og da snur tegnet: $x > -4$.
  \item $x \ge 4$ skrevet som $\langle 4, \rightarrow\rangle$ er feil. 4 er med, så det må være $[4, \rightarrow\rangle$.
  \item $4 < x + 2 \le 9$ gir $4 < x \le 9$ er feil. Du må trekke fra 2 i alle tre deler: $2 < x \le 7$.
\end{itemize}
\end{vanligfeil}

\nivaahode{1}

\oppg{1.1}{1}{\uh}
Løs ulikhetene.
\begin{tasks}(4)
  \task $x + 5 > 9$
  \task $3x \le 18$
  \task $\dfrac{x}{2} < -3$
  \task $x - 7 \ge -2$
\end{tasks}

\oppg{1.2}{1}{\uh}
Fullfør løsningen av $7 - 2x < 13$. Skriv tall og ulikhetstegn i rutene.

\medskip
Trekk fra 7: \quad $-2x$ \rute\ \rute

Del på $-2$: \quad $x$ \rute\ \rute

Intervall: \skrivlinje[3cm]

\oppg{1.3}{1}{\uh}
Løs ulikhetene. Pass på når du deler på et negativt tall.
\begin{tasks}(4)
  \task $-2x < 10$
  \task $-3x \ge 18$
  \task $4 - x > 1$
  \task $2 - 5x \le 12$
\end{tasks}

\oppg{1.4}{1}{\uh}
Skriv ulikheten og intervallet som hver tallinje viser.

\begin{center}
\begin{tabular}{@{}l l@{}}
a) & \begin{tikzpicture}[x=0.7cm,font=\small]
  \draw[-{Stealth[length=2mm]}] (-3.5,0) -- (6.5,0);
  \draw (-3,0.07) -- (-3,-0.07) node[below,font=\scriptsize,inner sep=2pt] {$-3$};
  \draw (-2,0.07) -- (-2,-0.07) node[below,font=\scriptsize,inner sep=2pt] {$-2$};
  \draw (-1,0.07) -- (-1,-0.07) node[below,font=\scriptsize,inner sep=2pt] {$-1$};
  \draw (0,0.07) -- (0,-0.07) node[below,font=\scriptsize,inner sep=2pt] {$0$};
  \draw (1,0.07) -- (1,-0.07) node[below,font=\scriptsize,inner sep=2pt] {$1$};
  \draw (2,0.07) -- (2,-0.07) node[below,font=\scriptsize,inner sep=2pt] {$2$};
  \draw (3,0.07) -- (3,-0.07) node[below,font=\scriptsize,inner sep=2pt] {$3$};
  \draw (4,0.07) -- (4,-0.07) node[below,font=\scriptsize,inner sep=2pt] {$4$};
  \draw (5,0.07) -- (5,-0.07) node[below,font=\scriptsize,inner sep=2pt] {$5$};
  \draw (6,0.07) -- (6,-0.07) node[below,font=\scriptsize,inner sep=2pt] {$6$};
  \draw[hoved,line width=2.2pt] (1.0,0) -- (6.5,0);
  \filldraw[hoved,fill=white,line width=0.9pt] (1.0,0) circle (2.6pt);
\end{tikzpicture}\\
b) & \begin{tikzpicture}[x=0.7cm,font=\small]
  \draw[-{Stealth[length=2mm]}] (-3.5,0) -- (6.5,0);
  \draw (-3,0.07) -- (-3,-0.07) node[below,font=\scriptsize,inner sep=2pt] {$-3$};
  \draw (-2,0.07) -- (-2,-0.07) node[below,font=\scriptsize,inner sep=2pt] {$-2$};
  \draw (-1,0.07) -- (-1,-0.07) node[below,font=\scriptsize,inner sep=2pt] {$-1$};
  \draw (0,0.07) -- (0,-0.07) node[below,font=\scriptsize,inner sep=2pt] {$0$};
  \draw (1,0.07) -- (1,-0.07) node[below,font=\scriptsize,inner sep=2pt] {$1$};
  \draw (2,0.07) -- (2,-0.07) node[below,font=\scriptsize,inner sep=2pt] {$2$};
  \draw (3,0.07) -- (3,-0.07) node[below,font=\scriptsize,inner sep=2pt] {$3$};
  \draw (4,0.07) -- (4,-0.07) node[below,font=\scriptsize,inner sep=2pt] {$4$};
  \draw (5,0.07) -- (5,-0.07) node[below,font=\scriptsize,inner sep=2pt] {$5$};
  \draw (6,0.07) -- (6,-0.07) node[below,font=\scriptsize,inner sep=2pt] {$6$};
  \draw[hoved,line width=2.2pt] (-3.5,0) -- (2.0,0);
  \filldraw[hoved,fill=hoved,line width=0.9pt] (2.0,0) circle (2.6pt);
\end{tikzpicture}\\
c) & \begin{tikzpicture}[x=0.7cm,font=\small]
  \draw[-{Stealth[length=2mm]}] (-3.5,0) -- (6.5,0);
  \draw (-3,0.07) -- (-3,-0.07) node[below,font=\scriptsize,inner sep=2pt] {$-3$};
  \draw (-2,0.07) -- (-2,-0.07) node[below,font=\scriptsize,inner sep=2pt] {$-2$};
  \draw (-1,0.07) -- (-1,-0.07) node[below,font=\scriptsize,inner sep=2pt] {$-1$};
  \draw (0,0.07) -- (0,-0.07) node[below,font=\scriptsize,inner sep=2pt] {$0$};
  \draw (1,0.07) -- (1,-0.07) node[below,font=\scriptsize,inner sep=2pt] {$1$};
  \draw (2,0.07) -- (2,-0.07) node[below,font=\scriptsize,inner sep=2pt] {$2$};
  \draw (3,0.07) -- (3,-0.07) node[below,font=\scriptsize,inner sep=2pt] {$3$};
  \draw (4,0.07) -- (4,-0.07) node[below,font=\scriptsize,inner sep=2pt] {$4$};
  \draw (5,0.07) -- (5,-0.07) node[below,font=\scriptsize,inner sep=2pt] {$5$};
  \draw (6,0.07) -- (6,-0.07) node[below,font=\scriptsize,inner sep=2pt] {$6$};
  \draw[hoved,line width=2.2pt] (0.0,0) -- (5.0,0);
  \filldraw[hoved,fill=hoved,line width=0.9pt] (0.0,0) circle (2.6pt);
  \filldraw[hoved,fill=white,line width=0.9pt] (5.0,0) circle (2.6pt);
\end{tikzpicture}\\
\end{tabular}
\end{center}

\nivaahode{2}

\oppg{1.5}{2}{\uh}
Løs ulikhetene. Svar som ulikhet og som intervall.
\begin{tasks}(2)
  \task $3x - 5 \le 2x + 4$
  \task $2(x - 3) > 5x + 3$
  \task $\dfrac{x + 2}{2} \ge \dfrac{x}{3}$
  \task $6 - 3(x + 1) < x - 1$
\end{tasks}

\oppg{1.6}{2}{\uh}
Løs de doble ulikhetene. Svar som intervall.
\begin{tasks}(2)
  \task $-4 \le 3x - 1 < 8$
  \task $1 < 5 - 2x \le 9$
\end{tasks}

\oppg{1.7}{2}{\uh}
Finn \textbf{den første feilen}. Løs riktig.
\[ 2(3 - x) + 4 \ge 5x - 11 \;\longrightarrow\; 6 - 2x + 4 \ge 5x - 11 \;\longrightarrow\; 10 - 2x \ge 5x - 11 \;\longrightarrow\; -7x \ge -21 \;\longrightarrow\; x \ge 3 \]

\oppg{1.8}{2}{\mh}
Abonnement A koster 99 kr i måneden og 15 kr per GB. Abonnement B koster 249 kr i måneden, uansett bruk.
\begin{tasks}(1)
  \task La $g$ være antall GB. Sett opp en ulikhet for når A er billigst, og løs den.
  \task Hvor mange hele GB kan du bruke i måneden hvis A skal være billigst? Svar i en setning.
\end{tasks}

\nivaahode{3}

\oppg{1.9}{3}{\uh}
Er påstanden alltid, noen ganger eller aldri sann? Begrunn svaret.
\begin{tasks}(1)
  \task Hvis $a < b$, så er $-a < -b$.
  \task Hvis $a < b$, så er $a^2 < b^2$.
  \task $2x + 1 > 2x - 3$ for alle tall $x$.
\end{tasks}

\oppg{1.10}{3}{\uh}
Løs $4 - 3x > 2x - 16$ på to måter.
Metode A: samle $x$-leddene på venstre side. Metode B: samle dem på høyre side.
Hvor snur du tegnet? Forklar hvorfor begge metodene gir samme svar.

\oppg{1.11}{3}{\uh}
Løs ulikheten $ax > 6$ for $x$. Hvordan avhenger svaret av om $a$ er positiv, negativ eller null? Forklar.

\begin{sjekk}
\begin{tasks}(2)
  \task Løs $2 - 4x \le 10$.
  \task Skriv $-2 \le x < 4$ som intervall.
  \task! Hvorfor snur vi ulikhetstegnet når vi deler på et negativt tall? Forklar med to tall.
\end{tasks}
\end{sjekk}

% ---------------------------------------------------------------------

\del{2}{Andregradsulikheter med fortegnslinje}

\begin{ordliste}
fortegnslinje & Figur som viser hvor et uttrykk er positivt (heltrukket) og negativt (stiplet). & \mittsprak\\
faktor & Et uttrykk som ganges med noe annet. I $(x+1)(x-3)$ er $x+1$ og $x-3$ faktorer. & \mittsprak\\
nullpunkt & En verdi av $x$ som gjør uttrykket lik 0. & \mittsprak\\
andregradsulikhet & En ulikhet der den høyeste potensen av $x$ er 2. & \mittsprak\\
union & $A \cup B$ betyr «A eller B»: tallene som er med i minst ett av intervallene. & \mittsprak\\
\end{ordliste}

\begin{mustkunne}
\textbf{Metode.} (1) Flytt alt over på én side, slik at den andre siden er 0. (2) Faktoriser (se 2.3).
(3) Tegn en fortegnslinje for hver faktor. (4) Tegn linja for produktet og les av der ulikheten er sann.

\textbf{Hvorfor virker det?} Et produkt kan bare skifte fortegn der en faktor er 0. To like fortegn gir pluss, to ulike gir minus.

\medskip
\begin{minipage}[t]{0.56\linewidth}
\vspace{0pt}
\begin{tikzpicture}[font=\small,fpos/.style={line width=1.1pt},fneg/.style={line width=1.1pt,dash pattern=on 4pt off 3pt}]
  \draw[graa,thin] (1.40,0.40) -- (1.40,-1.95);
  \draw[graa,thin] (3.30,0.40) -- (3.30,-1.95);
  \node[anchor=east] at (-0.15,-0.00) {$x + 1$};
  \draw[fneg] (0.00,-0.00) -- (1.40,-0.00);
  \draw[fpos] (1.40,-0.00) -- (3.30,-0.00);
  \draw[fpos] (3.30,-0.00) -- (4.70,-0.00);
  \node[fill=white,inner sep=1.5pt] at (1.40,-0.00) {$0$};
  \node[anchor=east] at (-0.15,-0.70) {$x - 3$};
  \draw[fneg] (0.00,-0.70) -- (1.40,-0.70);
  \draw[fneg] (1.40,-0.70) -- (3.30,-0.70);
  \draw[fpos] (3.30,-0.70) -- (4.70,-0.70);
  \node[fill=white,inner sep=1.5pt] at (3.30,-0.70) {$0$};
  \node[anchor=east] at (-0.15,-1.55) {$(x + 1) (x - 3)$};
  \draw[fpos,hoved] (0.00,-1.55) -- (1.40,-1.55);
  \draw[fneg,hoved] (1.40,-1.55) -- (3.30,-1.55);
  \draw[fpos,hoved] (3.30,-1.55) -- (4.70,-1.55);
  \node[fill=white,inner sep=1.5pt] at (1.40,-1.55) {$0$};
  \node[fill=white,inner sep=1.5pt] at (3.30,-1.55) {$0$};
  \node[below] at (1.40,-1.95) {$-1$};
  \node[below] at (3.30,-1.95) {$3$};
\end{tikzpicture}
\end{minipage}\hfill
\begin{minipage}[t]{0.42\linewidth}
\vspace{0pt}
\begin{tikzpicture}
\begin{axis}[width=6.2cm,height=4.6cm,axis lines=middle,xmin=-2,xmax=4.4,ymin=-5,ymax=5.5,
  xtick=\empty,ytick=\empty,xlabel={$x$},ylabel={$y$},font=\small,
  every axis x label/.style={at={(current axis.right of origin)},anchor=west},
  every axis y label/.style={at={(current axis.above origin)},anchor=south}]
  \addplot[domain=-1.8:3.8,samples=40,hoved,thick]{x^2-2*x-3};
  \addplot[domain=-1:3,samples=40,niv3,very thick]{x^2-2*x-3};
  \addplot[only marks,mark=*,mark size=1.8pt,mark options={fill=white}] coordinates {(-1,0) (3,0)};
  \node[anchor=north east] at (axis cs:-1,0) {$-1$};
  \node[anchor=north west] at (axis cs:3,0) {$3$};
\end{axis}
\end{tikzpicture}
\end{minipage}

\smallskip
$(x+1)(x-3) < 0$ der produktlinja er stiplet: $\langle -1, 3\rangle$. Grafen til $y = x^2 - 2x - 3$ ligger under $x$-aksen (rødt) i samme intervall.
\end{mustkunne}

\begin{eksempel}[1 — fortegnslinje]
Løs $x^2 - 2x - 3 < 0$.

\begin{stegtabell}
$(x+1)(x-3) < 0$ & Faktoriser: to tall med produkt $-3$ og sum $-2$ er $1$ og $-3$.\\
$x = -1$ og $x = 3$ & Nullpunktene. Tegn fortegnslinjene.\\
\end{stegtabell}

\smallskip
\begin{center}
\begin{tikzpicture}[font=\small,fpos/.style={line width=1.1pt},fneg/.style={line width=1.1pt,dash pattern=on 4pt off 3pt}]
  \draw[graa,thin] (1.40,0.40) -- (1.40,-1.95);
  \draw[graa,thin] (3.30,0.40) -- (3.30,-1.95);
  \node[anchor=east] at (-0.15,-0.00) {$x + 1$};
  \draw[fneg] (0.00,-0.00) -- (1.40,-0.00);
  \draw[fpos] (1.40,-0.00) -- (3.30,-0.00);
  \draw[fpos] (3.30,-0.00) -- (4.70,-0.00);
  \node[fill=white,inner sep=1.5pt] at (1.40,-0.00) {$0$};
  \node[anchor=east] at (-0.15,-0.70) {$x - 3$};
  \draw[fneg] (0.00,-0.70) -- (1.40,-0.70);
  \draw[fneg] (1.40,-0.70) -- (3.30,-0.70);
  \draw[fpos] (3.30,-0.70) -- (4.70,-0.70);
  \node[fill=white,inner sep=1.5pt] at (3.30,-0.70) {$0$};
  \node[anchor=east] at (-0.15,-1.55) {$(x + 1) (x - 3)$};
  \draw[fpos,hoved] (0.00,-1.55) -- (1.40,-1.55);
  \draw[fneg,hoved] (1.40,-1.55) -- (3.30,-1.55);
  \draw[fpos,hoved] (3.30,-1.55) -- (4.70,-1.55);
  \node[fill=white,inner sep=1.5pt] at (1.40,-1.55) {$0$};
  \node[fill=white,inner sep=1.5pt] at (3.30,-1.55) {$0$};
  \node[below] at (1.40,-1.95) {$-1$};
  \node[below] at (3.30,-1.95) {$3$};
\end{tikzpicture}
\end{center}

\begin{stegtabell}
$\langle -1, 3\rangle$ & Produktet er negativt (stiplet) mellom nullpunktene. Tegnet er $<$, så endepunktene er ikke med.\\
Kontroll: & $x = 0$ gir $-3 < 0$, som stemmer. $x = 4$ gir $16 - 8 - 3 = 5$, som ikke er negativt.\\
\end{stegtabell}
\end{eksempel}

\begin{eksempel}[2 — ikke del på $x$]
Løs $2x^2 \ge 8x$.

\begin{stegtabell}
$2x^2 - 8x \ge 0$ & Flytt alt over. Ikke del på $x$: vi vet ikke om $x$ er positiv eller negativ.\\
$2x(x - 4) \ge 0$ & Sett utenfor felles faktor $2x$. Faktoren 2 er positiv og endrer ikke fortegnet.\\
\end{stegtabell}

\smallskip
\begin{center}
\begin{tikzpicture}[font=\small,fpos/.style={line width=1.1pt},fneg/.style={line width=1.1pt,dash pattern=on 4pt off 3pt}]
  \draw[graa,thin] (1.40,0.40) -- (1.40,-1.95);
  \draw[graa,thin] (3.30,0.40) -- (3.30,-1.95);
  \node[anchor=east] at (-0.15,-0.00) {$x$};
  \draw[fneg] (0.00,-0.00) -- (1.40,-0.00);
  \draw[fpos] (1.40,-0.00) -- (3.30,-0.00);
  \draw[fpos] (3.30,-0.00) -- (4.70,-0.00);
  \node[fill=white,inner sep=1.5pt] at (1.40,-0.00) {$0$};
  \node[anchor=east] at (-0.15,-0.70) {$x - 4$};
  \draw[fneg] (0.00,-0.70) -- (1.40,-0.70);
  \draw[fneg] (1.40,-0.70) -- (3.30,-0.70);
  \draw[fpos] (3.30,-0.70) -- (4.70,-0.70);
  \node[fill=white,inner sep=1.5pt] at (3.30,-0.70) {$0$};
  \node[anchor=east] at (-0.15,-1.55) {$2x(x-4)$};
  \draw[fpos,hoved] (0.00,-1.55) -- (1.40,-1.55);
  \draw[fneg,hoved] (1.40,-1.55) -- (3.30,-1.55);
  \draw[fpos,hoved] (3.30,-1.55) -- (4.70,-1.55);
  \node[fill=white,inner sep=1.5pt] at (1.40,-1.55) {$0$};
  \node[fill=white,inner sep=1.5pt] at (3.30,-1.55) {$0$};
  \node[below] at (1.40,-1.95) {$0$};
  \node[below] at (3.30,-1.95) {$4$};
\end{tikzpicture}
\end{center}

\begin{stegtabell}
$x \le 0$ eller $x \ge 4$ & Produktet er positivt (heltrukket) utenfor nullpunktene. Tegnet er $\ge$, så endepunktene er med.\\
$\langle\leftarrow, 0] \cup [4, \rightarrow\rangle$ & To intervaller skrives med $\cup$ («eller»).\\
Kontroll: & $x = -1$ gir $2 \ge -8$, som stemmer. $x = 2$ gir $8 \ge 16$, som er usant.\\
\end{stegtabell}
\end{eksempel}

\begin{eksempel}[3 — når det ikke finnes nullpunkter]
Løs $x^2 - 4x + 5 > 0$.

\begin{stegtabell}
$D = (-4)^2 - 4 \cdot 1 \cdot 5 = -4$ & $D < 0$: ingen nullpunkter. Parabelen krysser ikke $x$-aksen.\\
$a = 1 > 0$ & Parabelen åpner oppover, så den ligger over $x$-aksen for alle $x$.\\
$\langle\leftarrow, \rightarrow\rangle$ & Svar: alle tall. Kontroll: $x^2 - 4x + 5 = (x-2)^2 + 1$. Et kvadrat er aldri negativt.\\
\end{stegtabell}
\end{eksempel}

\begin{vanligfeil}
\begin{itemize}[leftmargin=1.2em,itemsep=1pt,topsep=0pt]
  \item $x^2 > 4$, så $x > 2$ er feil. Også $x = -3$ virker. Riktig er $x < -2$ eller $x > 2$. Og $x^2 < 4$ gir $-2 < x < 2$.
  \item «$(x+1)(x-3) < 0$ betyr $x + 1 < 0$ eller $x - 3 < 0$» er feil. For $x = -2$ er begge faktorene negative, men produktet er $(-1)(-5) = 5 > 0$. Bruk fortegnslinje.
  \item $x^2 \ge 5x$, del på $x$, $x \ge 5$ er feil. Vi mister $x \le 0$. Riktig er $x(x - 5) \ge 0$, som gir $x \le 0$ eller $x \ge 5$.
\end{itemize}
\end{vanligfeil}

\nivaahode{1}

\oppg{2.1}{1}{\uh}
Vi skal løse $x^2 - 6x + 8 \le 0$. Fortegnslinjene til faktorene er tegnet. Produktlinja mangler.
\begin{tasks}(1)
  \task Vis at $(x-2)(x-4) = x^2 - 6x + 8$.
  \task Tegn av figuren og fyll inn produktlinja.
  \task Løs ulikheten. Svar som ulikhet og som intervall.
\end{tasks}

\begin{center}
\begin{tikzpicture}[font=\small,fpos/.style={line width=1.1pt},fneg/.style={line width=1.1pt,dash pattern=on 4pt off 3pt}]
  \draw[graa,thin] (1.40,0.40) -- (1.40,-1.95);
  \draw[graa,thin] (3.30,0.40) -- (3.30,-1.95);
  \node[anchor=east] at (-0.15,-0.00) {$x - 2$};
  \draw[fneg] (0.00,-0.00) -- (1.40,-0.00);
  \draw[fpos] (1.40,-0.00) -- (3.30,-0.00);
  \draw[fpos] (3.30,-0.00) -- (4.70,-0.00);
  \node[fill=white,inner sep=1.5pt] at (1.40,-0.00) {$0$};
  \node[anchor=east] at (-0.15,-0.70) {$x - 4$};
  \draw[fneg] (0.00,-0.70) -- (1.40,-0.70);
  \draw[fneg] (1.40,-0.70) -- (3.30,-0.70);
  \draw[fpos] (3.30,-0.70) -- (4.70,-0.70);
  \node[fill=white,inner sep=1.5pt] at (3.30,-0.70) {$0$};
  \node[anchor=east] at (-0.15,-1.55) {$(x - 2) (x - 4)$};
  \node[below] at (1.40,-1.95) {$2$};
  \node[below] at (3.30,-1.95) {$4$};
\end{tikzpicture}
\end{center}

\oppg{2.2}{1}{\uh}
Løs ulikhetene med fortegnslinje. Svar som intervall.
\begin{tasks}(3)
  \task $(x - 1)(x + 5) > 0$
  \task $(x + 2)(x - 5) \le 0$
  \task $x(x - 3) < 0$
\end{tasks}

\oppg{2.3}{1}{\uh}
Faktoriser og løs ulikhetene.
\begin{tasks}(3)
  \task $x^2 + 2x - 8 < 0$
  \task $x^2 - 7x + 10 \ge 0$
  \task $x^2 \le 9$
\end{tasks}

\nivaahode{2}

\oppg{2.4}{2}{\uh}
Flytt alt over på én side og løs ulikhetene.
\begin{tasks}(1)
  \task $x^2 > 3x + 10$
  \task $x^2 + 12 \le 7x$
  \task $2x^2 - 2x - 12 < 0$
\end{tasks}

\oppg{2.5}{2}{\uh}
Løs ulikhetene. Avgjør først om uttrykket har nullpunkter.
\begin{tasks}(2)
  \task $x^2 + 2x + 5 > 0$
  \task $(x + 3)^2 \le 0$
  \task $x^2 - 8x + 16 > 0$
  \task $x^2 + 6x + 10 \le 0$
\end{tasks}

\oppg{2.6}{2}{\uh}
Finn \textbf{den første feilen}. Løs riktig.
\[ x^2 - x > 12 \;\longrightarrow\; x^2 - x - 12 > 0 \;\longrightarrow\; (x - 3)(x + 4) > 0 \;\longrightarrow\; x < -4 \text{ eller } x > 3 \]

\oppg{2.7}{2}{\uh}
Figuren viser grafen til $f(x) = -x^2 + 2x + 8$.

\begin{center}
\begin{tikzpicture}
\begin{axis}[width=7.5cm,height=5.2cm,axis lines=middle,xmin=-3.3,xmax=5.3,ymin=-7.5,ymax=10.5,
  xtick={-3,-2,...,5},ytick={-6,-4,...,10},grid=major,grid style={graa!25},
  xlabel={$x$},ylabel={$f(x)$},font=\scriptsize,tick label style={font=\tiny},
  every axis x label/.style={at={(current axis.right of origin)},anchor=west},
  every axis y label/.style={at={(current axis.above origin)},anchor=south}]
  \addplot[domain=-2.8:4.8,samples=50,hoved,thick]{-x^2+2*x+8};
\end{axis}
\end{tikzpicture}
\end{center}
\begin{tasks}(1)
  \task Les av nullpunktene til $f$.
  \task Les av for hvilke $x$ som $f(x) \ge 0$, og for hvilke $x$ som $f(x) < 0$.
  \task Vis at $f(x) = -(x+2)(x-4)$. Bekreft svaret i b) med fortegnslinje. Faktoren $-1$ får en egen rad.
\end{tasks}

\nivaahode{3}

\oppg{2.8}{3}{\uh}
Er påstanden alltid, noen ganger eller aldri sann for et reelt tall $x$? Begrunn svaret.
\begin{tasks}(2)
  \task $x^2 > x$
  \task $(x - 3)^2 \ge 0$
  \task $x^2 + 1 < 0$
  \task $-x^2 \le 0$
\end{tasks}

\oppg{2.9}{3}{\uh}
Lag selv en andregradsulikhet med løsningsmengde $[-1, 6]$.
Lag også en med løsningsmengde «alle tall unntatt 3». Forklar hvorfor ulikhetene dine virker.

\oppg{2.10}{3}{\uh}
For hvilke verdier av $k$ er $x^2 + kx + 4 > 0$ for alle $x$? Begrunn svaret.
\hint{Tenk på grafen.}

\begin{sjekk}
\begin{tasks}(2)
  \task Løs $x^2 + x - 12 \ge 0$.
  \task Løs $-(x - 2)^2 < 0$.
  \task! Hvorfor kan vi ikke dele på $x$ når vi løser $x^2 \ge 7x$? Forklar.
\end{tasks}
\end{sjekk}
